\subsection{Обзор литературы по радиофотонике}

Радиофотонные РЛС опираются сразу на несколько направлений: физику и элементную
базу фотоники, методы СВЧ-обработки в оптической области, широкополосные
антенные решетки, фотонное аналого-цифровое преобразование и цифровую обработку
радиолокационных данных. Литературу по теме удобно рассматривать не по отдельным
компонентам, а по тому, какую роль они играют в полной радиолокационной системе.

Одной из наиболее близких к теме настоящей работы является статья
\cite{shumov2016concept}, в которой предложена концепция построения РЛС с активной
фазированной антенной решеткой на основе радиофотонных элементов. Авторы
рассматривают радиофотонику не как частную замену одного узла, а как системный
подход к построению станции: от синтезатора опорных частот и устройства формирования
зондирующего сигнала до оптического приемо-передающего модуля, волоконно-оптических
линий связи и оптического цифрового приемника. В статье также подчеркивается, что
основные выигрыши проявляются одновременно по нескольким направлениям: возрастает
разрешающая способность, снижаются потери в тракте передачи, уменьшается влияние
внешних электромагнитных воздействий, повышается стабильность характеристик и
сокращаются массогабаритные параметры антенной системы.

Существенная часть работы \cite{shumov2016concept} посвящена сравнению
традиционной и радиофотонной элементной базы. В качестве примера авторы анализируют
РЛС типа <<Воронеж-М>> и показывают, что переход к радиофотонным узлам может
существенно изменить расчетные тактико-технические характеристики системы. В
частности, расширение полосы до гигагерцового диапазона дает резкий рост разрешения
по дальности, а замена протяженных электрических линий на ВОЛС снижает суммарные
потери и массу антенной системы. Для настоящей работы этот результат важен как
обоснование общей инженерной мотивации: радиофотоника дает выигрыш не только в
абстрактных лабораторных параметрах, но и в характеристиках полной радиолокационной
системы.

Фундаментальные представления о фотонной элементной базе и возможностях
нанофотоники изложены в работе \cite{zaitsev2012nanophotonics}. Для радиолокационных
приложений здесь важны прежде всего методы управления оптическим излучением,
передачи сигналов по волокну, преобразования частоты и использования нелинейных
оптических эффектов. Эти идеи образуют физическую основу радиофотонных генераторов,
модуляторов, фильтров и оптических линий задержки. Однако для инженерной РЛС
недостаточно знать только свойства отдельных фотонных компонентов: необходимо
понимать, как они влияют на шумы, фазовую стабильность, динамический диапазон и
повторяемость измерений.

Обзор радиофотонных методов обработки аналоговых и цифровых СВЧ-сигналов
представлен в работе \cite{belousov2014radiophotonic}. Авторы рассматривают
применение радиофотоники в широком диапазоне длин волн --- от дециметрового до
миллиметрового, где классическая электронная обработка сталкивается с ограничениями
полосы и быстродействия. В контексте радиолокации особенно важны радиофотонные аналоговые
тракты, оптические фильтры, линии задержки и устройства предварительной обработки.
Такие узлы позволяют переносить часть операций до цифрового преобразования, что
уменьшает нагрузку на электронные АЦП и снижает требования к сверхскоростной
электронике.

Близкая по тематике работа \cite{gamilovskaya2015uwb} посвящена вариантам
реализации сверхширокополосных аналоговых радиофотонных процессоров. СШП-сигналы
являются естественным объектом интереса для РЛС, поскольку именно большая полоса
повышает разрешающую способность по дальности. При этом сверхширокополосность
создает ряд трудностей: требуется сохранить амплитудно-фазовую линейность тракта,
избежать частотно-зависимых искажений и обеспечить достаточно высокий динамический
диапазон. Радиофотонные процессоры рассматриваются как один из способов обработки
таких сигналов без полной передачи всех требований на электронную часть приемника.

Особое место в литературе занимает фотонное формирование и распределение сигналов
для антенных решеток. В классической АФАР каждый излучатель или группа излучателей
требует управляемого фазового сдвига, а при широкой полосе --- управляемой временной
задержки. Работа \cite{lee1995photonicwideband} демонстрирует, что фотонные методы
могут применяться для широкополосных антенных решеток и диаграммообразования.
Идея состоит в том, что оптическая линия задержки реализует физическую задержку
сигнала, а не только фазовый поворот на одной частоте. Поэтому фотонное
диаграммообразование особенно привлекательно для систем, где требуется сохранить
форму сверхширокополосного импульса или обеспечить согласованную работу во всем
диапазоне частот.

Современный обзор радиофотоники для радиолокационных систем приведен в
\cite{pan2014microwave}. В этой работе радиофотонные методы рассматриваются в
контексте генерации стабильных СВЧ-сигналов, широкополосной передачи, обработки
сигналов, формирования диаграммы направленности и аналого-цифрового преобразования.
Важная мысль данного обзора состоит в том, что фотонные методы способны расширить
возможности РЛС именно за счет сочетания большой полосы, малых потерь в линиях,
устойчивости к электромагнитным воздействиям и совместимости с распределенными
антенными структурами. Этот вывод согласуется с концепцией, предложенной в
\cite{shumov2016concept}, но рассматривает проблему в более широком международном
контексте.

Одним из наиболее сложных узлов радиофотонной РЛС остается приемник. Если
система использует сверхширокополосный сигнал, то электронный АЦП должен обеспечить
одновременно высокую частоту дискретизации, достаточное эффективное число бит,
малый джиттер и широкий аналоговый входной диапазон. Эти требования трудно
совместить в одном серийном устройстве. Поэтому в литературе активно исследуются
фотонные АЦП и схемы оптической предобработки. Обзор таких подходов приведен в
\cite{starikov2015photonicadc}; фундаментальные ограничения и практические аспекты
фотонного аналого-цифрового преобразования рассмотрены в \cite{valley2007photonicadc}.

В работе \cite{khilo2012photonicadc} фотонный АЦП рассматривается как способ
уменьшить влияние электронного джиттера, который становится одним из основных
ограничителей при оцифровке высокочастотных сигналов. Для радиолокации это имеет
прямое значение: фазовая ошибка, вызванная нестабильностью временной привязки,
может привести к ошибке оценки дальности и ухудшению когерентного накопления.
Тем не менее литература показывает, что фотонные АЦП пока не являются универсальной
заменой электронным: остаются ограничения по сложности, стоимости, числу бит,
линейности и доступности серийных решений. Поэтому в прикладных системах часто
выбирается компромиссный подход, при котором оптическая часть выполняет генерацию,
передачу или предварительную обработку, а окончательное выделение информативных
компонент производится электронным синхронным детектором и цифровым трактом.

Экспериментальные демонстрации радиофотонных радаров подтверждают реализуемость
таких схем. В проекте PHODIR рассматривалась фотонно-цифровая радиолокационная
система, где фотонные методы использовались для формирования и обработки
радиолокационных сигналов \cite{phodir2016}. В литературе также описываются
оптические широкополосные радары с перестраиваемой полосой, а также отдельные
экспериментальные узлы: волоконно-оптические линии передачи СВЧ-сигналов,
фотодиодные модули, оптические генераторы и линии задержки. В совокупности эти
работы показывают, что радиофотоника уже вышла за пределы чисто теоретической
концепции, но построение устойчивой измерительной системы по-прежнему требует
решения задач согласования, калибровки и управления.

Для настоящей работы важен и другой пласт литературы --- методы построения
радиолокационных изображений по частотным измерениям. В георадиолокации широко
используются импульсные, FMCW- и SFCW-схемы. Обзор современных методов обработки
данных GPR представлен в \cite{sun2019advanced}, а особенности SFCW-систем и их
практические преимущества рассматриваются в \cite{eide2019advanced}. В SFCW-радаре
широкая полоса синтезируется не за счет излучения короткого импульса, а за счет
последовательного измерения комплексного отклика на наборе частот. Такой подход
хорошо согласуется с задачами синхронного детектирования: на каждой частоте система
может стабильно измерить амплитуду и фазу отклика, после чего временная картина
восстанавливается с помощью обратного преобразования Фурье \cite{noon1996stepped}.

Связь между радиофотонной аппаратной частью и SFCW-методом состоит в том, что
оба подхода требуют строгой синхронизации и хорошей фазовой стабильности. Если
частотная перестройка, опорный сигнал синхронного детектора и регистрация данных
несогласованы, то комплексный S-параметр теряет физический смысл: фаза начинает
содержать не только информацию о задержке отраженной волны, но и ошибки аппаратного
тракта. Поэтому система управления радиофотонным радаром является не вспомогательным
программным блоком, а частью измерительной методики. Она должна задавать частотную
сетку, управлять последовательностью измерений, обеспечивать повторяемую запись
данных и поддерживать процедуры калибровки.

Значительная часть публикаций сосредоточена либо на отдельных элементах
радиофотоники, либо на концептуальных схемах полной РЛС. Промежуточному
инженерному уровню уделяется меньше внимания: как организовать управление
лабораторным радиофотонным радаром, синхронизировать перестройку частоты и
детектирование, сохранить комплексный отклик в пригодной для обработки форме
и проверить работоспособность системы на экспериментальных данных. Именно этот
уровень определяет прикладную направленность настоящей работы.

Таким образом, статья \cite{shumov2016concept} задает системную мотивацию перехода
к радиофотонной РЛС, работы \cite{zaitsev2012nanophotonics,
belousov2014radiophotonic, gamilovskaya2015uwb, pan2014microwave} раскрывают
физические и схемотехнические основы радиофотонной обработки, публикации
\cite{starikov2015photonicadc, valley2007photonicadc, khilo2012photonicadc}
описывают ограничения приемно-цифрового тракта, а исследования
\cite{sun2019advanced, eide2019advanced, noon1996stepped} связывают аппаратные
возможности с методами построения радиолокационных изображений. На этом основании
можно сформулировать задачу настоящей работы как разработку управляемой
измерительной системы, в которой преимущества радиофотонной аппаратуры используются
для получения и обработки комплексного радиолокационного отклика.
